% CS 422 Assignment 3 report. Build with:  python3 report/make_report.py
% (writes report/generated/*.tex from results/<run>/ and runs pdflatex twice).
\documentclass[11pt]{article}
\usepackage[T1]{fontenc}
\usepackage{lmodern}
\usepackage[utf8]{inputenc}
\usepackage[margin=1in]{geometry}
\usepackage{microtype}
\usepackage{graphicx}
\usepackage{booktabs}
\usepackage{array}
\usepackage{longtable}
\usepackage{float}
\usepackage{caption}
\usepackage{xcolor}
\usepackage{listings}
\usepackage[normalem]{ulem}
\usepackage{amsmath}
\usepackage[colorlinks=true,linkcolor=blue!50!black,urlcolor=blue!60!black,citecolor=blue!50!black]{hyperref}

% generated by report/make_report.py -- do not edit
\newcommand{\plotdir}{../results/20261018-133538/plots/}
\newcommand{\rundir}{results/20261018-133538}
\newcommand{\repourl}{https://github.com/philwisniewski/cs-422-assignment-3}
\newcommand{\reporef}{main}
\expandafter\def\csname meta@kernel\endcsname{6.17.0-1022-nvidia}
\expandafter\def\csname meta@uname\endcsname{Linux jensen 6.17.0-1022-nvidia \#22-Ubuntu SMP PREEMPT\_DYNAMIC Wed May 27 19:31:58 UTC 2026 aarch64 aarch64}
\expandafter\def\csname meta@python\endcsname{3.12.3}
\expandafter\def\csname meta@availcc\endcsname{reno cubic bbr gentle\_reno}
\expandafter\def\csname meta@defaultcc\endcsname{cubic}
\expandafter\def\csname meta@command\endcsname{python3 run\_experiments.py --servers servers.csv -n 10 -t 10 -r 5 --seed 42}
\expandafter\def\csname meta@created\endcsname{2026-10-18 13:35:38 -0400}
\expandafter\def\csname meta@finished\endcsname{2026-10-18 14:22:02 -0400}
\expandafter\def\csname meta@repdest\endcsname{lon.speedtest.clouvider.net:5207}
\expandafter\def\csname meta@seed\endcsname{42}
\expandafter\def\csname meta@numdests\endcsname{10}
\expandafter\def\csname meta@duration\endcsname{10}
\expandafter\def\csname meta@interval\endcsname{0.2}
\expandafter\def\csname meta@reps\endcsname{5}
\expandafter\def\csname meta@pause\endcsname{2}
\expandafter\def\csname meta@blksize\endcsname{131072}
\expandafter\def\csname meta@serversparsed\endcsname{189}
\expandafter\def\csname meta@selattempts\endcsname{15}
\expandafter\def\csname cccheck@cubic\endcsname{cubic}
\expandafter\def\csname cccheck@reno\endcsname{reno}
\expandafter\def\csname cccheck@bbr\endcsname{bbr}
\expandafter\def\csname cccheck@gentle_reno\endcsname{gentle\_reno}
\expandafter\def\csname bpf@stat_init\endcsname{47}
\expandafter\def\csname bpf@stat_loss_events\endcsname{15}
\expandafter\def\csname bpf@stat_ss_exit_delay\endcsname{2}
\expandafter\def\csname bpf@stat_ss_exit_bw\endcsname{44}
\expandafter\def\csname bpf@stat_undo\endcsname{5}
\expandafter\def\csname order@total\endcsname{200}
\expandafter\def\csname order@ok\endcsname{190}
\expandafter\def\csname order@failed\endcsname{10}
\expandafter\def\csname order@faileddests\endcsname{\texttt{speed1.\allowbreak{}fiberby.\allowbreak{}dk:9212}}
\expandafter\def\csname agg@cubic@algorithm\endcsname{cubic}
\expandafter\def\csname agg@cubic@algorithm_actual\endcsname{cubic}
\expandafter\def\csname agg@cubic@runs\endcsname{48.00}
\expandafter\def\csname agg@cubic@destinations\endcsname{10.00}
\expandafter\def\csname agg@cubic@samples\endcsname{2400.00}
\expandafter\def\csname agg@cubic@goodput_mean_mbps\endcsname{132.17}
\expandafter\def\csname agg@cubic@goodput_median_mbps\endcsname{76.68}
\expandafter\def\csname agg@cubic@goodput_p5_mbps\endcsname{0.12}
\expandafter\def\csname agg@cubic@goodput_p95_mbps\endcsname{394.43}
\expandafter\def\csname agg@cubic@srtt_mean_ms\endcsname{137.02}
\expandafter\def\csname agg@cubic@srtt_median_ms\endcsname{130.62}
\expandafter\def\csname agg@cubic@srtt_p95_ms\endcsname{250.75}
\expandafter\def\csname agg@cubic@rtt_inflation_mean_ms\endcsname{7.29}
\expandafter\def\csname agg@cubic@retrans_rate_mean_pct\endcsname{1.63}
\expandafter\def\csname agg@cubic@retrans_per_run_mean\endcsname{1120.06}
\expandafter\def\csname agg@cubic@rto_per_run_mean\endcsname{0.02}
\expandafter\def\csname agg@reno@algorithm\endcsname{reno}
\expandafter\def\csname agg@reno@algorithm_actual\endcsname{reno}
\expandafter\def\csname agg@reno@runs\endcsname{47.00}
\expandafter\def\csname agg@reno@destinations\endcsname{10.00}
\expandafter\def\csname agg@reno@samples\endcsname{2350.00}
\expandafter\def\csname agg@reno@goodput_mean_mbps\endcsname{104.02}
\expandafter\def\csname agg@reno@goodput_median_mbps\endcsname{53.28}
\expandafter\def\csname agg@reno@goodput_p5_mbps\endcsname{0.12}
\expandafter\def\csname agg@reno@goodput_p95_mbps\endcsname{322.68}
\expandafter\def\csname agg@reno@srtt_mean_ms\endcsname{139.16}
\expandafter\def\csname agg@reno@srtt_median_ms\endcsname{132.86}
\expandafter\def\csname agg@reno@srtt_p95_ms\endcsname{253.10}
\expandafter\def\csname agg@reno@rtt_inflation_mean_ms\endcsname{7.63}
\expandafter\def\csname agg@reno@retrans_rate_mean_pct\endcsname{1.86}
\expandafter\def\csname agg@reno@retrans_per_run_mean\endcsname{1044.00}
\expandafter\def\csname agg@reno@rto_per_run_mean\endcsname{0.00}
\expandafter\def\csname agg@bbr@algorithm\endcsname{bbr}
\expandafter\def\csname agg@bbr@algorithm_actual\endcsname{bbr}
\expandafter\def\csname agg@bbr@runs\endcsname{48.00}
\expandafter\def\csname agg@bbr@destinations\endcsname{10.00}
\expandafter\def\csname agg@bbr@samples\endcsname{2400.00}
\expandafter\def\csname agg@bbr@goodput_mean_mbps\endcsname{171.19}
\expandafter\def\csname agg@bbr@goodput_median_mbps\endcsname{108.42}
\expandafter\def\csname agg@bbr@goodput_p5_mbps\endcsname{0.12}
\expandafter\def\csname agg@bbr@goodput_p95_mbps\endcsname{633.61}
\expandafter\def\csname agg@bbr@srtt_mean_ms\endcsname{146.47}
\expandafter\def\csname agg@bbr@srtt_median_ms\endcsname{135.80}
\expandafter\def\csname agg@bbr@srtt_p95_ms\endcsname{270.14}
\expandafter\def\csname agg@bbr@rtt_inflation_mean_ms\endcsname{14.84}
\expandafter\def\csname agg@bbr@retrans_rate_mean_pct\endcsname{5.09}
\expandafter\def\csname agg@bbr@retrans_per_run_mean\endcsname{7759.90}
\expandafter\def\csname agg@bbr@rto_per_run_mean\endcsname{0.04}
\expandafter\def\csname agg@gentle_reno@algorithm\endcsname{gentle\_reno}
\expandafter\def\csname agg@gentle_reno@algorithm_actual\endcsname{gentle\_reno}
\expandafter\def\csname agg@gentle_reno@runs\endcsname{47.00}
\expandafter\def\csname agg@gentle_reno@destinations\endcsname{10.00}
\expandafter\def\csname agg@gentle_reno@samples\endcsname{2350.00}
\expandafter\def\csname agg@gentle_reno@goodput_mean_mbps\endcsname{25.52}
\expandafter\def\csname agg@gentle_reno@goodput_median_mbps\endcsname{10.10}
\expandafter\def\csname agg@gentle_reno@goodput_p5_mbps\endcsname{1.27}
\expandafter\def\csname agg@gentle_reno@goodput_p95_mbps\endcsname{74.82}
\expandafter\def\csname agg@gentle_reno@srtt_mean_ms\endcsname{135.61}
\expandafter\def\csname agg@gentle_reno@srtt_median_ms\endcsname{130.62}
\expandafter\def\csname agg@gentle_reno@srtt_p95_ms\endcsname{247.84}
\expandafter\def\csname agg@gentle_reno@rtt_inflation_mean_ms\endcsname{4.46}
\expandafter\def\csname agg@gentle_reno@retrans_rate_mean_pct\endcsname{0.01}
\expandafter\def\csname agg@gentle_reno@retrans_per_run_mean\endcsname{16.55}
\expandafter\def\csname agg@gentle_reno@rto_per_run_mean\endcsname{0.00}
\expandafter\def\csname cwnd@cubic@seg\endcsname{1071}
\expandafter\def\csname cwnd@cubic@pathbdp\endcsname{47}
\expandafter\def\csname cwnd@reno@seg\endcsname{832}
\expandafter\def\csname cwnd@reno@pathbdp\endcsname{37}
\expandafter\def\csname cwnd@bbr@seg\endcsname{3289}
\expandafter\def\csname cwnd@bbr@pathbdp\endcsname{146}
\expandafter\def\csname cwnd@gentle_reno@seg\endcsname{149}
\expandafter\def\csname cwnd@gentle_reno@pathbdp\endcsname{7}
\newcommand{\stat}[2]{\csname #1@#2\endcsname}
\newcommand{\agg}[2]{\csname agg@#1@#2\endcsname}
\newcommand{\meta}[1]{\csname meta@#1\endcsname}
\newcommand{\bpfstat}[1]{\csname bpf@#1\endcsname}
\newcommand{\grParams}{\href{https://github.com/philwisniewski/cs-422-assignment-3/blob/main/bpf/gentle_reno.bpf.c\#L27-L38}{\texttt{bpf/gentle\_reno.bpf.c:27-38}}}
\newcommand{\grState}{\href{https://github.com/philwisniewski/cs-422-assignment-3/blob/main/bpf/gentle_reno.bpf.c\#L41-L50}{\texttt{bpf/gentle\_reno.bpf.c:41-50}}}
\newcommand{\grInit}{\href{https://github.com/philwisniewski/cs-422-assignment-3/blob/main/bpf/gentle_reno.bpf.c\#L102-L116}{\texttt{bpf/gentle\_reno.bpf.c:102-116}}}
\newcommand{\grSsthresh}{\href{https://github.com/philwisniewski/cs-422-assignment-3/blob/main/bpf/gentle_reno.bpf.c\#L121-L129}{\texttt{bpf/gentle\_reno.bpf.c:121-129}}}
\newcommand{\grCongAvoid}{\href{https://github.com/philwisniewski/cs-422-assignment-3/blob/main/bpf/gentle_reno.bpf.c\#L132-L147}{\texttt{bpf/gentle\_reno.bpf.c:132-147}}}
\newcommand{\grPktsAcked}{\href{https://github.com/philwisniewski/cs-422-assignment-3/blob/main/bpf/gentle_reno.bpf.c\#L150-L202}{\texttt{bpf/gentle\_reno.bpf.c:150-202}}}
\newcommand{\grSetState}{\href{https://github.com/philwisniewski/cs-422-assignment-3/blob/main/bpf/gentle_reno.bpf.c\#L205-L215}{\texttt{bpf/gentle\_reno.bpf.c:205-215}}}
\newcommand{\grUndo}{\href{https://github.com/philwisniewski/cs-422-assignment-3/blob/main/bpf/gentle_reno.bpf.c\#L218-L225}{\texttt{bpf/gentle\_reno.bpf.c:218-225}}}
\newcommand{\grOps}{\href{https://github.com/philwisniewski/cs-422-assignment-3/blob/main/bpf/gentle_reno.bpf.c\#L228-L236}{\texttt{bpf/gentle\_reno.bpf.c:228-236}}}
\newcommand{\pyConnect}{\href{https://github.com/philwisniewski/cs-422-assignment-3/blob/main/iperf3_client.py\#L164-L191}{\texttt{iperf3\_client.py:164-191}}}
\newcommand{\pySample}{\href{https://github.com/philwisniewski/cs-422-assignment-3/blob/main/iperf3_client.py\#L275-L332}{\texttt{iperf3\_client.py:275-332}}}
\newcommand{\pyRun}{\href{https://github.com/philwisniewski/cs-422-assignment-3/blob/main/iperf3_client.py\#L447-L572}{\texttt{iperf3\_client.py:447-572}}}
\newcommand{\pySetCC}{\href{https://github.com/philwisniewski/cs-422-assignment-3/blob/main/tcp_info.py\#L138-L140}{\texttt{tcp\_info.py:138-140}}}
\newcommand{\pyGetCC}{\href{https://github.com/philwisniewski/cs-422-assignment-3/blob/main/tcp_info.py\#L132-L135}{\texttt{tcp\_info.py:132-135}}}
\newcommand{\pyPreflight}{\href{https://github.com/philwisniewski/cs-422-assignment-3/blob/main/run_experiments.py\#L105-L126}{\texttt{run\_experiments.py:105-126}}}
\newcommand{\pyBpfSnapshot}{\href{https://github.com/philwisniewski/cs-422-assignment-3/blob/main/run_experiments.py\#L79-L102}{\texttt{run\_experiments.py:79-102}}}
\newcommand{\pySelect}{\href{https://github.com/philwisniewski/cs-422-assignment-3/blob/main/run_experiments.py\#L150-L198}{\texttt{run\_experiments.py:150-198}}}
\newcommand{\pyRunExperiments}{\href{https://github.com/philwisniewski/cs-422-assignment-3/blob/main/run_experiments.py\#L213-L267}{\texttt{run\_experiments.py:213-267}}}
\newcommand{\pyRunStats}{\href{https://github.com/philwisniewski/cs-422-assignment-3/blob/main/analysis.py\#L163-L200}{\texttt{analysis.py:163-200}}}
\newcommand{\pyBuildTables}{\href{https://github.com/philwisniewski/cs-422-assignment-3/blob/main/analysis.py\#L203-L294}{\texttt{analysis.py:203-294}}}
\newcommand{\pyChooseRep}{\href{https://github.com/philwisniewski/cs-422-assignment-3/blob/main/analysis.py\#L314-L352}{\texttt{analysis.py:314-352}}}
\newcommand{\pyPlotTs}{\href{https://github.com/philwisniewski/cs-422-assignment-3/blob/main/analysis.py\#L435-L446}{\texttt{analysis.py:435-446}}}
\newcommand{\pyPlotScatter}{\href{https://github.com/philwisniewski/cs-422-assignment-3/blob/main/analysis.py\#L471-L486}{\texttt{analysis.py:471-486}}}
\newcommand{\pyPlotRegimes}{\href{https://github.com/philwisniewski/cs-422-assignment-3/blob/main/analysis.py\#L536-L574}{\texttt{analysis.py:536-574}}}
\newcommand{\pyMakePlots}{\href{https://github.com/philwisniewski/cs-422-assignment-3/blob/main/analysis.py\#L577-L687}{\texttt{analysis.py:577-687}}}
\newcommand{\pyReport}{\href{https://github.com/philwisniewski/cs-422-assignment-3/blob/main/report/make_report.py\#L337-L376}{\texttt{report/make\_report.py:337-376}}}
\expandafter\def\csname file@bpfgentlerenobpfc\endcsname{\href{https://github.com/philwisniewski/cs-422-assignment-3/blob/main/bpf/gentle_reno.bpf.c}{\texttt{bpf/gentle\_reno.bpf.c}}}
\expandafter\def\csname file@bpfMakefile\endcsname{\href{https://github.com/philwisniewski/cs-422-assignment-3/blob/main/bpf/Makefile}{\texttt{bpf/Makefile}}}
\expandafter\def\csname file@runallsh\endcsname{\href{https://github.com/philwisniewski/cs-422-assignment-3/blob/main/run_all.sh}{\texttt{run\_all.sh}}}
\expandafter\def\csname file@runexperimentspy\endcsname{\href{https://github.com/philwisniewski/cs-422-assignment-3/blob/main/run_experiments.py}{\texttt{run\_experiments.py}}}
\expandafter\def\csname file@iperfclientpy\endcsname{\href{https://github.com/philwisniewski/cs-422-assignment-3/blob/main/iperf3_client.py}{\texttt{iperf3\_client.py}}}
\expandafter\def\csname file@tcpinfopy\endcsname{\href{https://github.com/philwisniewski/cs-422-assignment-3/blob/main/tcp_info.py}{\texttt{tcp\_info.py}}}
\expandafter\def\csname file@analysispy\endcsname{\href{https://github.com/philwisniewski/cs-422-assignment-3/blob/main/analysis.py}{\texttt{analysis.py}}}
\expandafter\def\csname file@serverspy\endcsname{\href{https://github.com/philwisniewski/cs-422-assignment-3/blob/main/servers.py}{\texttt{servers.py}}}
\expandafter\def\csname file@toolslocaltestbedsh\endcsname{\href{https://github.com/philwisniewski/cs-422-assignment-3/blob/main/tools/local_testbed.sh}{\texttt{tools/local\_testbed.sh}}}
\expandafter\def\csname file@docsapseudocodetxt\endcsname{\href{https://github.com/philwisniewski/cs-422-assignment-3/blob/main/docs/a2_pseudocode.txt}{\texttt{docs/a2\_pseudocode.txt}}}
\expandafter\def\csname file@docsimplementationmd\endcsname{\href{https://github.com/philwisniewski/cs-422-assignment-3/blob/main/docs/implementation.md}{\texttt{docs/implementation.md}}}
\expandafter\def\csname file@reportmakereportpy\endcsname{\href{https://github.com/philwisniewski/cs-422-assignment-3/blob/main/report/make_report.py}{\texttt{report/make\_report.py}}}
\expandafter\def\csname file@reportreporttex\endcsname{\href{https://github.com/philwisniewski/cs-422-assignment-3/blob/main/report/report.tex}{\texttt{report/report.tex}}}
\newcommand{\file}[1]{\csname file@#1\endcsname}


\lstset{basicstyle=\ttfamily\small, columns=fullflexible, keepspaces=true,
        frame=single, rulecolor=\color{black!25}, backgroundcolor=\color{black!3},
        xleftmargin=4pt, xrightmargin=4pt, aboveskip=6pt, belowskip=6pt,
        breaklines=true, showstringspaces=false,
        % listings reads bytes, so map the pseudocode's UTF-8 symbols explicitly
        literate={β}{{$\beta$}}1 {←}{{$\leftarrow$}}1 {∞}{{$\infty$}}1
                 {·}{{$\cdot$}}1 {−}{{$-$}}1 {≥}{{$\geq$}}1}
\captionsetup{font=small, labelfont=bf}
\setlength{\parskip}{4pt}
\setlength{\emergencystretch}{3em}
\sloppy
\newcommand{\alg}[1]{\texttt{#1}}
\newcommand{\GR}{Gentle Reno}

\newcommand{\aiuse}{%
AI tools (Claude, via Claude Code) were used to help brainstorm the design, draft
the eBPF program, the experiment and analysis scripts, the local test network, the
documentation, and the generator for this report. We reviewed the generated code,
built and ran it on our machine, and checked the behaviour of the algorithm on live
connections (loss cut ratio, slow-start exit counters, and the algorithm reported by
the socket).}
% ---------------------------------------------------------------------------

\title{\textbf{CS 422 Assignment 3: TCP Congestion Control Implementation}\\[4pt]
\large Option 1: \GR{} as an eBPF \texttt{struct\_ops} congestion control}
\author{Daniel Chindris \and Saverio Miller \and Richie Tan \and Philip Wisniewski \and Emma Zheng}
\date{CS 422: Computer Networks, Fall 2026\\[2pt]
Repository: \href{\repourl}{\nolinkurl{github.com/philwisniewski/cs-422-assignment-3}}}

\begin{document}
\maketitle

\begin{abstract}
\noindent We implemented the congestion control algorithm from Assignment~2, Question~3(b),
``\GR{}'', as an eBPF \texttt{struct\_ops} program that registers a kernel
\texttt{tcp\_congestion\_ops} named \alg{gentle\_reno}. Our Assignment~2 Python iperf3 client selects it on the data
socket with \texttt{setsockopt(TCP\_CONGESTION)} and reads back the algorithm actually used.
We compared it with CUBIC, Reno and BBR on \meta{numdests} public iperf3 servers, using
\meta{reps} repetitions of \meta{duration}\,s tests sampled every \meta{interval}\,s
(Linux \meta{kernel}). \GR{} had the lowest RTT inflation (\agg{gentle_reno}{rtt_inflation_mean_ms}\,ms vs.\
\agg{cubic}{rtt_inflation_mean_ms}\,ms for CUBIC and \agg{bbr}{rtt_inflation_mean_ms}\,ms for BBR) and
almost no retransmissions (\agg{gentle_reno}{retrans_rate_mean_pct}\,\% of bytes). It also had much lower
throughput (mean \agg{gentle_reno}{goodput_mean_mbps}\,Mbit/s vs.\ \agg{cubic}{goodput_mean_mbps},
\agg{reno}{goodput_mean_mbps} and \agg{bbr}{goodput_mean_mbps}\,Mbit/s). The cause is that its
slow-start exit fires early, after which its $+0.5$\,packet/RTT growth cannot fill a
long-RTT path within 10\,s.
\end{abstract}

\tableofcontents
\newpage

% ===========================================================================
\section{Algorithm (Assignment 2, Q3(b))}
\label{sec:design}

\GR{} keeps Reno's loss-driven AIMD structure but makes both halves gentler, and adds
RTT- and throughput-based rules for leaving slow start without a loss:
\begin{itemize}
  \item \textbf{Loss:} multiplicative decrease by $\beta = 0.7$ instead of $0.5$, at most once per loss episode.
  \item \textbf{Growth:} additive increase of $a = 0.5$ packet per RTT instead of $1$.
  \item \textbf{RTT:} leave slow start (without a cut) once the RTT rises more than $D = \max(4\,\text{ms}, \text{min\_rtt}/8)$ above the minimum, which means a queue is forming.
  \item \textbf{Throughput:} leave slow start once the delivery rate has grown by less than 25\,\% for three rounds in a row, which means the pipe is full.
\end{itemize}
The original pseudocode (\file{docsapseudocodetxt}) is:
\lstinputlisting{../docs/a2_pseudocode.txt}

% ===========================================================================
\section{Implementation}
\label{sec:impl}

The program is \file{bpfgentlerenobpfc}. It is compiled with clang to a BPF object and
registered with \texttt{bpftool struct\_ops register}, after which \alg{gentle\_reno} appears in
\texttt{/proc/sys/net/ipv4/tcp\_available\_congestion\_control} next to \alg{cubic}, \alg{reno}
and \alg{bbr}. The struct that registers the hooks is at \grOps.

\subsection{Division of work between the kernel and the program}
The Linux TCP stack still does everything except the window policy: ACK processing, loss
detection (dup ACKs, SACK, RACK, RTO timers), fast recovery with PRR (RFC~6937),
retransmission, and RTT estimation (srtt, rttvar, RTO). It calls our hooks at each decision point:

\begin{center}\small
\begin{tabular}{>{\raggedright\arraybackslash}p{2.6cm}>{\raggedright\arraybackslash}p{4.0cm}>{\raggedright\arraybackslash}p{5.3cm}l}
\toprule
Hook & Called when & What \GR{} does & Signal \\
\midrule
\texttt{gr\_init} & connection established & reset per-connection state, start the first round & -- \\
\texttt{gr\_ssthresh} & once per loss episode (enter Recovery after 3 dup ACKs/SACK/RACK, or Loss after RTO) & return $\max(\lfloor 7\,\text{cwnd}/10\rfloor, 2)$ & loss \\
\texttt{gr\_cong\_avoid} & each ACK that advances \texttt{snd\_una}, outside recovery & slow start: $+1$ per ACKed packet; else $+1$ packet per $2\cdot\text{cwnd}$ ACKed ($=+0.5$/RTT) & RTT (ACK clock) \\
\texttt{gr\_pkts\_acked} & each ACK, with its RTT sample & track min RTT; at each round end evaluate the slow-start exit rules & RTT, throughput \\
\texttt{gr\_set\_state} & CA state change & after an RTO re-arm the slow-start exit & loss (timeout) \\
\texttt{gr\_undo\_cwnd} & kernel proves a loss spurious (DSACK, timestamps) & restore $\max(\text{cwnd}, \text{prior\_cwnd})$, like Reno & loss \\
\bottomrule
\end{tabular}
\end{center}

\subsection{How the congestion window is updated}

\paragraph{Loss (\grSsthresh).}
\texttt{gr\_ssthresh} returns $0.7\cdot\text{cwnd}$ (at least 2). In fast recovery, PRR moves
cwnd smoothly down to that value, so after recovery $\text{cwnd} = \text{ssthresh} = 0.7\cdot\text{cwnd}_{\text{before}}$.
Because the kernel calls \texttt{ssthresh} only when it \emph{enters} recovery and not once per lost
packet, the cut happens at most once per loss episode (about once per RTT). On an RTO the kernel sets
$\text{cwnd}=1$ after calling \texttt{ssthresh}, so slow start restarts toward $0.7\cdot\text{cwnd}_{\text{before}}$,
as in the pseudocode. The loss signals we log are retransmitted segments per interval
($\Delta$\texttt{tcpi\_total\_retrans}), RTO count (\texttt{tcpi\_total\_rto}) and \texttt{ca\_state}
(3 = Recovery after dup ACKs, 4 = Loss after an RTO).
On a shallow-buffer test path we measured \texttt{ssthresh}/cwnd-before-loss at 0.70--0.73 for
\alg{gentle\_reno}, against 0.50--0.55 for \alg{reno} and about 0.7 for \alg{cubic}.

\paragraph{Growth and RTT (\grCongAvoid).}
Growth happens per round trip. \texttt{tcp\_cong\_avoid\_ai(tp, 2*cwnd, acked)} adds one packet after
$2\cdot\text{cwnd}$ packets are ACKed. A full window of ACKs arrives once per RTT, so this is $+0.5$
packet/RTT. The window grows only while the sender is cwnd-limited (\texttt{tcp\_is\_cwnd\_limited} logic), as in
every kernel algorithm.

\paragraph{RTT and throughput in slow start (\grPktsAcked).}
On every ACK the program tracks \texttt{min\_rtt\_us} (connection minimum) and
\texttt{round\_min\_rtt\_us} (minimum within the current round). A \emph{round} starts by recording
$\text{end\_seq}=\texttt{snd\_nxt}$ and ends when \texttt{snd\_una} passes it, i.e.\ one RTT. This is
the round definition that CUBIC's HyStart uses. At the end of each slow-start round:
\begin{itemize}
  \item \textbf{Throughput:} $\text{bw} = (\texttt{delivered} - \texttt{delivered}_{\text{round start}})\cdot 10^6 / \Delta t_{\mu s}$
  in packets/s. If $4\,\text{bw} \ge 5\,\text{full\_bw}$, then $\text{full\_bw}\leftarrow\text{bw}$ and the flat counter resets.
  Otherwise \texttt{flat\_rounds} is incremented. Slow start ends after 3 flat rounds.
  \item \textbf{RTT:} slow start ends if $\text{round\_min\_rtt} > \text{min\_rtt} + \max(4\,\text{ms}, \text{min\_rtt}/8)$.
\end{itemize}
Leaving slow start sets $\text{ssthresh} = \text{cwnd}$ with no cut, so the window reached is kept and
growth continues at $+0.5$/RTT.

\subsection{Initial state and parameters}
Parameters are \texttt{\#define}s at \grParams; per-connection state (32 bytes in
\texttt{icsk\_ca\_priv}, within the kernel's 104-byte limit) is at \grState; initialisation is
\grInit. All arithmetic is integer because BPF has no floating point.

\begin{center}\small
\begin{tabular}{l>{\raggedright\arraybackslash}p{5.6cm}>{\raggedright\arraybackslash}p{4.9cm}}
\toprule
Quantity & Value & Where \\
\midrule
cwnd & 10 packets (kernel IW10) & kernel default \\
ssthresh & $\infty$ (\texttt{TCP\_INFINITE\_SSTHRESH}): start in slow start & kernel default \\
min\_rtt, round\_min\_rtt & $\infty$ (\texttt{0xffffffff}) & \texttt{gr\_init} \\
full\_bw, flat\_rounds, ss\_done & 0 & \texttt{gr\_init} \\
round & \texttt{snd\_nxt}, \texttt{tcp\_mstamp}, \texttt{delivered} at init & \texttt{start\_round} \\
$\beta$ & 7/10 & \texttt{BETA\_NUM}/\texttt{BETA\_DEN} \\
$a$ & 1/2 packet per RTT & \texttt{AI\_DIVISOR = 2} \\
Minimum cwnd after a cut & 2 & \texttt{MIN\_CWND} \\
$D$ & $\max(4\,\text{ms}, \text{min\_rtt}/8)$ & \texttt{SS\_DELAY\_MIN\_US}, \texttt{SS\_DELAY\_SHIFT} \\
Throughput plateau & $<1.25\times$ growth for 3 rounds & \texttt{SS\_BW\_GROWTH\_*}, \texttt{SS\_FLAT\_ROUNDS} \\
\bottomrule
\end{tabular}
\end{center}

The program also keeps global counters (\texttt{stat\_init}, \texttt{stat\_loss\_events},
\texttt{stat\_ss\_exit\_delay}, \texttt{stat\_ss\_exit\_bw}, \texttt{stat\_undo}). They are read with
\texttt{sudo make -C bpf stats} and snapshotted before and after each experiment into
\texttt{metadata.json} (\pyBpfSnapshot).

\subsection{Changes from the Assignment 2 pseudocode}
\begin{enumerate}
  \item \textbf{The RTT rule uses the minimum RTT sample in the round, not srtt.} srtt is a smoothed
  average that lags several RTTs behind and is inflated by delayed ACKs and jitter. The per-round
  minimum reacts within one round and ignores single noisy samples. Linux HyStart makes the same choice. $D$ is unchanged.
  \item \textbf{The slow-start exit is evaluated at most once per slow-start phase} (\texttt{ss\_done})
  and is re-armed after an RTO (\texttt{gr\_set\_state}, \grSetState). The pseudocode left this implicit.
  \item \textbf{Cut and recovery mechanics come from the kernel.} \texttt{ssthresh()} supplies only the
  target, and PRR moves cwnd to it during recovery. The end state matches the pseudocode's ``cwnd $\leftarrow$ ssthresh''.
  \item \textbf{Undo of spurious cuts.} The kernel requires an \texttt{undo\_cwnd} hook (\grUndo).
  It restores the pre-loss window when a retransmission turns out to be unnecessary, as Reno does.
  \item \textbf{Growth only when cwnd-limited.} This prevents window inflation when the application or the
  receive window is the real limit.
  \item \textbf{Fixed-point parameters:} $\beta = 7/10$ in integers, and $a = 0.5$ as ``$+1$ per $2\cdot$cwnd ACKs''.
\end{enumerate}

% ===========================================================================
\section{Selecting the algorithm on the data socket}
\label{sec:select}

The Assignment~2 client (\file{iperfclientpy}) sets the algorithm on the \emph{data} socket before
\texttt{connect()}, so the whole connection, including slow start, uses it (\pyConnect, which calls
\pySetCC):
\begin{lstlisting}[language=Python]
sock.setsockopt(socket.IPPROTO_TCP, socket.TCP_CONGESTION, b"gentle_reno")
sock.connect(addr)
actual = sock.getsockopt(socket.IPPROTO_TCP, socket.TCP_CONGESTION, 16)
\end{lstlisting}
After the connection is up the client reads the algorithm back (\pyGetCC) and stores it as
\texttt{algorithm\_actual} in every sample row and every table. In every run of this experiment the requested and actual
algorithm matched (the ``actual'' column of Table~\ref{tab:agg} and the panel titles of the figures).
Before any test, \texttt{preflight} (\pyPreflight) opens a socket for each algorithm and records what the kernel
reports: \alg{cubic}$\to$\alg{\csname cccheck@cubic\endcsname},
\alg{reno}$\to$\alg{\csname cccheck@reno\endcsname}, \alg{bbr}$\to$\alg{\csname cccheck@bbr\endcsname},
\alg{gentle\_reno}$\to$\alg{\csname cccheck@gentle_reno\endcsname}.

% ===========================================================================
\section{Experimental setup}
\label{sec:setup}

\begin{center}\small
\begin{tabular}{>{\raggedright\arraybackslash}p{4cm}>{\raggedright\arraybackslash}p{11.5cm}}
\toprule
Host & \texttt{\meta{uname}} \\
Kernel & \texttt{\meta{kernel}} (aarch64, Ubuntu) \\
Available algorithms & \texttt{\meta{availcc}} (default \texttt{\meta{defaultcc}}) \\
Python & \meta{python} \\
Command & \texttt{\meta{command}} \\
Run directory & \texttt{\rundir} (started \meta{created}, finished \meta{finished}) \\
Destinations & \meta{numdests}, drawn at random (seed \meta{seed}) from \meta{serversparsed} servers in \texttt{servers.csv}; candidates that failed a 3\,s probe were replaced (\meta{selattempts} attempts, \pySelect) \\
Duration / sampling & \meta{duration}\,s per test, sampled every \meta{interval}\,s, \meta{pause}\,s pause between tests \\
Repetitions & \meta{reps} of every destination $\times$ algorithm: \stat{order}{total} scheduled tests, \stat{order}{ok} completed \\
Failed tests & \stat{order}{failed}, all at \stat{order}{faileddests} (\texttt{connect\_error}: the server refused the connection). They are reported, not retried into the statistics. \\
\bottomrule
\end{tabular}
\end{center}

\noindent\textbf{Destinations.} All four algorithms used the same \meta{numdests} destinations:

\begin{center}\small\begin{tabular}{rlllr}
\toprule
\# & Destination & Location & Provider & Gbit/s \\
\midrule
1 & \texttt{41.\allowbreak{}110.\allowbreak{}39.\allowbreak{}130:5201} & Algiers, DZ & DATAPACKET & 10 \\
2 & \texttt{akl.\allowbreak{}linetest.\allowbreak{}nz:5301} & Auckland, NZ & 2Degrees & 10 \\
3 & \texttt{lon.\allowbreak{}speedtest.\allowbreak{}clouvider.\allowbreak{}net:5207} & London, GB & Cloudvider & 10 \\
4 & \texttt{speedtest.\allowbreak{}sin1.\allowbreak{}sg.\allowbreak{}leaseweb.\allowbreak{}net:5207} & Singapore, SG & LeaseWeb & 10 \\
5 & \texttt{speedtest.\allowbreak{}wdc2.\allowbreak{}us.\allowbreak{}leaseweb.\allowbreak{}net:5203} & Washington, US & LeaseWeb & 10 \\
6 & \texttt{72.\allowbreak{}1.\allowbreak{}86.\allowbreak{}46:30006} & Auckland, NZ & FortiSASE & 10 \\
7 & \texttt{speed1.\allowbreak{}fiberby.\allowbreak{}dk:9212} & Copenhagen, DK & Fiberby & 25 \\
8 & \texttt{41.\allowbreak{}210.\allowbreak{}185.\allowbreak{}162:5201} & Kampala, UG & DATAPACKET & 2x10 \\
9 & \texttt{speedtest.\allowbreak{}nyc1.\allowbreak{}us.\allowbreak{}leaseweb.\allowbreak{}net:5207} & New York City, US & LeaseWeb & 10 \\
10 & \texttt{speedtest.\allowbreak{}phx1.\allowbreak{}us.\allowbreak{}leaseweb.\allowbreak{}net:5206} & Phoenix, US & LeaseWeb & 10 \\
\bottomrule
\end{tabular}
\end{center}

\noindent\textbf{Test order.} For each repetition, for each destination, the four algorithms run
back to back. Their order is rotated by $(\text{rep} + \text{destination index}) \bmod 4$, so no
algorithm always goes first and slow drift in path conditions is spread across algorithms
(\pyRunExperiments). The full order is listed in Appendix~\ref{app:order} and in
\texttt{test\_order.csv}.

\noindent\textbf{Logged per sample} (\pySample; \texttt{samples/*.csv}): goodput
($\Delta$\texttt{tcpi\_bytes\_acked}$/\Delta t$), \texttt{snd\_cwnd}, \texttt{snd\_ssthresh},
srtt, rttvar, min\_rtt, RTO, kernel delivery rate, and the loss signals: $\Delta$\texttt{tcpi\_total\_retrans}
per interval, \texttt{bytes\_retrans}, \texttt{lost}, \texttt{total\_rto} and \texttt{ca\_state}.

\noindent\textbf{Representative destination.} \texttt{\meta{repdest}}. The analysis picks the destination
with complete data for every algorithm whose median goodput is closest to the median over all
destinations (\pyChooseRep). All four algorithms are plotted for this destination with
identical axis limits.

% ===========================================================================
\section{Results}
\label{sec:results}

\subsection{Time series and scatter plots (A2 Q2(b), all four algorithms)}

Figures~\ref{fig:ts} and~\ref{fig:sc} repeat the Assignment~2 Q2(b) plots for all four algorithms on
the representative destination, side by side with shared axes. The individual per-algorithm
versions are in Appendix~\ref{app:plots}. They are generated by \texttt{make\_plots} (\pyMakePlots),
using \texttt{plot\_timeseries} (\pyPlotTs) and \texttt{plot\_scatter} (\pyPlotScatter).

\begin{figure}[H]
  \centering
  \includegraphics[width=\textwidth]{\plotdir q3_timeseries_compare.pdf}
  \caption{Time series at \texttt{\meta{repdest}}, repetition 1: congestion window and ssthresh,
  RTT (srtt, srtt$\pm$rttvar, min\_rtt), retransmitted segments per interval (loss signal, symlog),
  and goodput. Columns are CUBIC, Reno, BBR and \GR{}, each titled with the requested and actual algorithm;
  every row uses the same y-axis across algorithms. File: \texttt{plots/q3\_timeseries\_compare.pdf}.}
  \label{fig:ts}
\end{figure}

\begin{figure}[H]
  \centering
  \includegraphics[width=\textwidth]{\plotdir q3_scatter_compare.pdf}
  \caption{Goodput against cwnd, srtt and retransmissions per interval at \texttt{\meta{repdest}},
  all repetitions, with the Pearson correlation in each panel. Same axes for all algorithms.
  File: \texttt{plots/q3\_scatter\_compare.pdf}.}
  \label{fig:sc}
\end{figure}

\subsection{Aggregate throughput and RTT statistics}

Table~\ref{tab:agg} pools every \meta{interval}\,s interval sample of every completed run (all destinations
and repetitions). Percentiles use \texttt{numpy.percentile} with its default linear interpolation
between order statistics (Hyndman--Fan type~7). RTT inflation, retransmission rate and RTO count are
computed per run and then averaged. The code is \texttt{build\_tables} (\pyBuildTables) and
\texttt{run\_stats} (\pyRunStats).

\begin{table}[H]
  \centering\small
  \begin{tabular}{lrrrr}
\toprule
Statistic & \textbf{CUBIC} & \textbf{Reno} & \textbf{BBR} & \textbf{Gentle Reno} \\
\midrule
Runs / interval samples & 48 / 2400 & 47 / 2350 & 48 / 2400 & 47 / 2350 \\
Goodput mean [Mbit/s] & 132.2 & 104.0 & 171.2 & 25.5 \\
Goodput median [Mbit/s] & 76.7 & 53.3 & 108.4 & 10.1 \\
Goodput p5 [Mbit/s] & 0.12 & 0.12 & 0.12 & 1.27 \\
Goodput p95 [Mbit/s] & 394.4 & 322.7 & 633.6 & 74.8 \\
srtt mean [ms] & 137.0 & 139.2 & 146.5 & 135.6 \\
srtt median [ms] & 130.6 & 132.9 & 135.8 & 130.6 \\
srtt p95 [ms] & 250.8 & 253.1 & 270.1 & 247.8 \\
RTT inflation (srtt $-$ min\_rtt) [ms] & 7.29 & 7.63 & 14.84 & 4.46 \\
Retransmitted bytes [\%] & 1.63 & 1.86 & 5.09 & 0.01 \\
Retransmitted segments / run & 1120 & 1044 & 7760 & 17 \\
RTOs / run & 0.02 & 0.00 & 0.04 & 0.00 \\
\bottomrule
\end{tabular}

  \caption{Aggregate statistics per algorithm. Source: \texttt{tables/aggregate\_by\_algorithm.csv}.
  The ``actual'' algorithm recorded on every socket was
  \alg{\agg{cubic}{algorithm_actual}}, \alg{\agg{reno}{algorithm_actual}}, \alg{\agg{bbr}{algorithm_actual}},
  \alg{\agg{gentle_reno}{algorithm_actual}} respectively.}
  \label{tab:agg}
\end{table}

Tables~\ref{tab:pdg} and~\ref{tab:pdr} break the results down by destination. Each cell is the mean
$\pm$ standard deviation over repetitions of the per-run mean. Bold marks the best algorithm on that
path, and a superscript gives the number of completed runs when it is not \meta{reps}.

\begin{table}[H]
  \centering\small
  \begin{tabular}{>{\raggedright\arraybackslash}p{5.4cm}rrrr}
\toprule
Destination & \textbf{CUBIC} & \textbf{Reno} & \textbf{BBR} & \textbf{Gentle Reno} \\
\midrule
\texttt{41.\allowbreak{}110.\allowbreak{}39.\allowbreak{}130:5201} & 92.5 $\pm$ 10.2 & 53.7 $\pm$ 19.4 & \textbf{95.6 $\pm$ 9.2} & 13.0 $\pm$ 10.6 \\
\texttt{akl.\allowbreak{}linetest.\allowbreak{}nz:5301} & 67.0 $\pm$ 7.3 & 36.8 $\pm$ 1.0 & \textbf{72.9 $\pm$ 10.2} & 9.4 $\pm$ 9.0 \\
\texttt{lon.\allowbreak{}speedtest.\allowbreak{}clouvider.\allowbreak{}net:5207} & 65.5 $\pm$ 46.7 & 77.5 $\pm$ 22.7 & \textbf{164.0 $\pm$ 13.9} & 11.1 $\pm$ 0.8 \\
\texttt{speedtest.\allowbreak{}sin1.\allowbreak{}sg.\allowbreak{}leaseweb.\allowbreak{}net:5207} & 28.2 $\pm$ 13.4 & 34.2 $\pm$ 10.4 & \textbf{56.5 $\pm$ 7.7} & 3.9 $\pm$ 2.3 \\
\texttt{speedtest.\allowbreak{}wdc2.\allowbreak{}us.\allowbreak{}leaseweb.\allowbreak{}net:5203} & 320.8 $\pm$ 84.9 & 234.3 $\pm$ 15.7 & \textbf{414.3 $\pm$ 83.9} & 53.0 $\pm$ 11.1 \\
\texttt{72.\allowbreak{}1.\allowbreak{}86.\allowbreak{}46:30006} & 48.4 $\pm$ 4.7 & 37.1 $\pm$ 1.2 & \textbf{74.3 $\pm$ 1.0} & 8.1 $\pm$ 4.8 \\
\texttt{speed1.\allowbreak{}fiberby.\allowbreak{}dk:9212} & 6.6 $\pm$ 0.3$^{(3)}$ & 6.3 $\pm$ 0.3$^{(2)}$ & 6.4 $\pm$ 0.1$^{(3)}$ & \textbf{22.3 $\pm$ 13.4$^{(2)}$} \\
\texttt{41.\allowbreak{}210.\allowbreak{}185.\allowbreak{}162:5201} & 44.6 $\pm$ 5.9 & 35.5 $\pm$ 9.3 & \textbf{60.1 $\pm$ 3.8} & 3.7 $\pm$ 0.6 \\
\texttt{speedtest.\allowbreak{}nyc1.\allowbreak{}us.\allowbreak{}leaseweb.\allowbreak{}net:5207} & 373.9 $\pm$ 68.8 & 275.4 $\pm$ 44.2 & \textbf{445.1 $\pm$ 192.4} & 100.1 $\pm$ 97.6 \\
\texttt{speedtest.\allowbreak{}phx1.\allowbreak{}us.\allowbreak{}leaseweb.\allowbreak{}net:5206} & 224.2 $\pm$ 49.7 & 190.8 $\pm$ 23.3 & \textbf{256.7 $\pm$ 52.5} & 28.6 $\pm$ 3.7 \\
\midrule
\textit{All destinations} & 132.2 $\pm$ 131.0 & 104.0 $\pm$ 94.1 & \textbf{171.2 $\pm$ 163.1} & 25.5 $\pm$ 44.0 \\
\bottomrule
\end{tabular}

  \caption{Mean goodput per destination [Mbit/s]. Source: \texttt{tables/algorithm\_comparison.csv}.}
  \label{tab:pdg}
\end{table}

\begin{table}[H]
  \centering\small
  \begin{tabular}{>{\raggedright\arraybackslash}p{5.4cm}rrrr}
\toprule
Destination & \textbf{CUBIC} & \textbf{Reno} & \textbf{BBR} & \textbf{Gentle Reno} \\
\midrule
\texttt{41.\allowbreak{}110.\allowbreak{}39.\allowbreak{}130:5201} & 136.0 $\pm$ 5.1 & 148.0 $\pm$ 21.7 & 144.4 $\pm$ 3.6 & \textbf{132.3 $\pm$ 1.7} \\
\texttt{akl.\allowbreak{}linetest.\allowbreak{}nz:5301} & 190.1 $\pm$ 2.3 & 190.1 $\pm$ 1.6 & 199.9 $\pm$ 6.2 & \textbf{188.4 $\pm$ 2.0} \\
\texttt{lon.\allowbreak{}speedtest.\allowbreak{}clouvider.\allowbreak{}net:5207} & 102.8 $\pm$ 6.4 & \textbf{102.7 $\pm$ 6.3} & 108.4 $\pm$ 3.6 & 103.1 $\pm$ 3.4 \\
\texttt{speedtest.\allowbreak{}sin1.\allowbreak{}sg.\allowbreak{}leaseweb.\allowbreak{}net:5207} & 241.8 $\pm$ 6.9 & 239.7 $\pm$ 10.3 & 260.8 $\pm$ 12.8 & \textbf{238.5 $\pm$ 10.6} \\
\texttt{speedtest.\allowbreak{}wdc2.\allowbreak{}us.\allowbreak{}leaseweb.\allowbreak{}net:5203} & 35.1 $\pm$ 1.8 & \textbf{33.0 $\pm$ 1.9} & 45.1 $\pm$ 6.9 & 37.2 $\pm$ 3.1 \\
\texttt{72.\allowbreak{}1.\allowbreak{}86.\allowbreak{}46:30006} & 200.7 $\pm$ 5.6 & 203.8 $\pm$ 4.7 & 209.3 $\pm$ 3.9 & \textbf{196.9 $\pm$ 6.2} \\
\texttt{speed1.\allowbreak{}fiberby.\allowbreak{}dk:9212} & 121.2 $\pm$ 5.3$^{(3)}$ & 118.4 $\pm$ 0.2$^{(2)}$ & 127.7 $\pm$ 10.2$^{(3)}$ & \textbf{113.1 $\pm$ 1.2$^{(2)}$} \\
\texttt{41.\allowbreak{}210.\allowbreak{}185.\allowbreak{}162:5201} & 251.6 $\pm$ 2.8 & 254.5 $\pm$ 5.4 & 256.0 $\pm$ 3.4 & \textbf{246.1 $\pm$ 2.1} \\
\texttt{speedtest.\allowbreak{}nyc1.\allowbreak{}us.\allowbreak{}leaseweb.\allowbreak{}net:5207} & \textbf{30.3 $\pm$ 1.9} & 30.8 $\pm$ 1.4 & 41.3 $\pm$ 17.1 & 32.1 $\pm$ 2.2 \\
\texttt{speedtest.\allowbreak{}phx1.\allowbreak{}us.\allowbreak{}leaseweb.\allowbreak{}net:5206} & \textbf{54.2 $\pm$ 2.9} & 58.2 $\pm$ 2.6 & 64.4 $\pm$ 3.9 & 54.8 $\pm$ 1.8 \\
\midrule
\textit{All destinations} & 137.0 $\pm$ 79.8 & 139.2 $\pm$ 81.2 & 146.5 $\pm$ 80.5 & \textbf{135.6 $\pm$ 78.4} \\
\bottomrule
\end{tabular}

  \caption{Mean srtt per destination [ms]; bold = lowest. Source: \texttt{tables/algorithm\_comparison.csv}.}
  \label{tab:pdr}
\end{table}

\begin{figure}[H]
  \centering
  \includegraphics[width=0.92\textwidth]{\plotdir q3_regimes.pdf}
  \caption{Per-run mean goodput and srtt for each destination (top), and goodput against the path's
  min RTT and against the retransmission rate (bottom). One marker per run. Generated by
  \texttt{plot\_regimes} (\pyPlotRegimes). File: \texttt{plots/q3\_regimes.pdf}.}
  \label{fig:regimes}
\end{figure}

\subsection{What the eBPF program saw}
Counters read from the program after the experiment: \bpfstat{stat_init} connections used
\alg{gentle\_reno}. Slow start ended through the throughput-plateau rule \bpfstat{stat_ss_exit_bw} times
and through the RTT-rise rule \bpfstat{stat_ss_exit_delay} times. There were \bpfstat{stat_loss_events}
loss episodes (\texttt{ssthresh} calls) and \bpfstat{stat_undo} spurious-loss undos.
The mean congestion window was \stat{cwnd}{gentle_reno@seg} segments for \GR{},
compared with \stat{cwnd}{cubic@seg} (CUBIC), \stat{cwnd}{reno@seg} (Reno) and
\stat{cwnd}{bbr@seg} (BBR). As a fraction of each path's BDP (the largest
delivery-rate $\times$ min-RTT estimate any algorithm reached on that destination), the mean
window was \stat{cwnd}{gentle_reno@pathbdp}\,\% for \GR{}, \stat{cwnd}{cubic@pathbdp}\,\% for CUBIC,
\stat{cwnd}{reno@pathbdp}\,\% for Reno and \stat{cwnd}{bbr@pathbdp}\,\% for BBR.

% ===========================================================================
\section{Discussion of the observed differences}
\label{sec:discussion}

\paragraph{BBR: highest throughput, most queueing and loss.}
BBR had the highest mean goodput on \emph{every} destination except the anomalous
\texttt{speed1.fiberby.dk} (Table~\ref{tab:pdg}): \agg{bbr}{goodput_mean_mbps}\,Mbit/s overall.
It also had the highest RTT inflation (\agg{bbr}{rtt_inflation_mean_ms}\,ms) and by far the most
retransmissions (\agg{bbr}{retrans_rate_mean_pct}\,\% of bytes, \agg{bbr}{retrans_per_run_mean} segments per run).
BBR (v1 in this kernel) does not treat loss as a congestion signal. It paces at its bottleneck-bandwidth
estimate, probes with $1.25\times$ gain, and caps cwnd at about $2\times$ the estimated BDP. Its window
sat above the path BDP (\stat{cwnd}{bbr@pathbdp}\,\%), so it keeps a standing queue and overflows
shallow buffers. Its throughput does not collapse after loss the way the loss-based
algorithms' does. Figure~\ref{fig:ts} shows this: BBR's cwnd stays near 4700 segments despite repeated
retransmission bursts, while CUBIC and Reno drop to a few hundred segments after their first losses
around $t\approx1$--$2$\,s.

\paragraph{CUBIC vs.\ Reno.}
Both are loss-based and saw similar RTT (\agg{cubic}{srtt_mean_ms} vs.\ \agg{reno}{srtt_mean_ms}\,ms) and
retransmission rates. CUBIC was faster on average (\agg{cubic}{goodput_mean_mbps} vs.\
\agg{reno}{goodput_mean_mbps}\,Mbit/s). It cuts by only 30\,\% (Reno halves), and its cubic growth
function depends on time since the last loss, not on RTT. On the long-RTT paths here (most min RTTs
are 100--250\,ms), Reno's $+1$ packet/RTT takes many seconds to regain the window after a cut. CUBIC
returns to its previous maximum in a few seconds. The gap is largest on the low-RTT, high-bandwidth
US paths (\texttt{wdc2}, \texttt{nyc1}), where the window can grow the most within 10\,s.

\paragraph{\GR{}: lowest delay and loss, lowest throughput.}
\GR{} had the lowest mean srtt (\agg{gentle_reno}{srtt_mean_ms}\,ms), the lowest RTT inflation
(\agg{gentle_reno}{rtt_inflation_mean_ms}\,ms) and essentially no retransmissions
(\agg{gentle_reno}{retrans_per_run_mean} segments per run, against more than 1000 for the others). Its
mean goodput was only \agg{gentle_reno}{goodput_mean_mbps}\,Mbit/s. The data point to one cause:
\begin{itemize}
  \item The slow-start exit fired in nearly every connection, almost always through the
  throughput-plateau rule (\bpfstat{stat_ss_exit_bw} of \bpfstat{stat_init}), at a window of roughly
  50--200 segments in most runs. That is far below the path BDP: the mean window was only
  \stat{cwnd}{gentle_reno@pathbdp}\,\% of BDP. In Figure~\ref{fig:ts} the \GR{} cwnd flattens within
  the first second and never approaches the 4000+ segments the other algorithms reach.
  \item After the exit, growth is $+0.5$ packet/RTT. At 100\,ms RTT that is only about $+50$ packets in
  a 10\,s test, so the window never reaches the buffer and almost never causes a loss. This explains the
  low RTT and loss, and also the low throughput. The scatter plot (Figure~\ref{fig:sc}) shows a high
  goodput--cwnd correlation for \GR{}: it is window-limited by its own cwnd, not by the network.
  \item The one run whose slow start did not exit early (\texttt{nyc1}, repetition~1) reached
  292\,Mbit/s, matching CUBIC on that path. The control law can use the bandwidth once the window
  is large.
\end{itemize}
The likely reason the plateau rule fires early is that some early rounds are not cwnd-limited.
During the first RTTs the receiver's advertised window (still being autotuned up) or ACK
aggregation can hold delivery flat for a few rounds. The rule then counts those rounds as ``pipe
full''. BBR's full-bandwidth detector avoids this by ignoring application- or receiver-limited
rounds. Adding the same check (skip rounds in which the flow was not cwnd-limited) and requiring the
throughput plateau to agree with a modest RTT rise would make the exit much less premature. The
gentle $a=0.5$ increase also makes recovery from an early exit very slow on long-RTT paths. A
HyStart-like fallback, or a CUBIC-style growth curve after an early exit, would address that. The
algorithm as designed in Assignment~2 is very conservative. It almost never builds a queue, but
in 10\,s tests over 30--250\,ms paths it leaves most of the bandwidth unused. Over a much longer
transfer its throughput would keep rising toward the path capacity.

\paragraph{The \texttt{speed1.fiberby.dk} anomaly.}
On this path every loss-based and BBR run had a median of 0.12\,Mbit/s with a short burst, and half
of the scheduled tests failed (\stat{order}{failed} failures, all on this server). This points to
server-side throttling. \GR{} looks ``best'' there only because 2 of its runs completed with a short
continuous transfer. We do not draw conclusions from this destination.

% ===========================================================================
\section{Reproduction}
\label{sec:repro}

\subsection{Kernel and required features}
\begin{itemize}
  \item Kernel used for all results: \textbf{\texttt{\meta{kernel}}} (Ubuntu, aarch64); exact
  \texttt{uname -a} in Section~\ref{sec:setup} and in \texttt{metadata.json}. Any Linux $\ge$ 5.13 with BTF works.
  \item Kernel configuration:\\
  \hspace*{1em}\texttt{CONFIG\_BPF\_SYSCALL=y}, \texttt{CONFIG\_BPF\_JIT=y}\\
  \hspace*{1em}\texttt{CONFIG\_DEBUG\_INFO\_BTF=y} (\texttt{/sys/kernel/btf/vmlinux} must exist)\\
  \hspace*{1em}\texttt{CONFIG\_TCP\_CONG\_ADVANCED=y}, and \texttt{CONFIG\_TCP\_CONG\_BBR} for the BBR comparison\\
  The program uses the BPF \texttt{tcp\_congestion\_ops} \texttt{struct\_ops} and the
  \texttt{tcp\_slow\_start} / \texttt{tcp\_cong\_avoid\_ai} kfuncs.
  \item Tools: clang $\ge$ 12, libbpf headers (\texttt{libbpf-dev}), bpftool
  (\texttt{linux-tools-\$(uname -r)} or \texttt{linux-tools-generic}), Python 3 with numpy and matplotlib.
  \item Root (or \texttt{CAP\_BPF} + \texttt{CAP\_NET\_ADMIN}) to register the program and select
  \alg{bbr}/\alg{gentle\_reno}.
\end{itemize}

\subsection{Build, load and check}
\begin{lstlisting}[language=bash]
sudo apt install clang llvm libbpf-dev linux-tools-common linux-tools-$(uname -r) \
                 python3-numpy python3-matplotlib
make -C bpf             # dump vmlinux.h from the running kernel, clang -target bpf
sudo make -C bpf load   # bpftool struct_ops register gentle_reno.bpf.o
cat /proc/sys/net/ipv4/tcp_available_congestion_control  # ... gentle_reno
sudo make -C bpf status
sudo make -C bpf stats  # program counters
python3 -c "import socket; s=socket.socket(); \
  s.setsockopt(6, 13, b'gentle_reno'); print(s.getsockopt(6, 13, 16))"
sudo make -C bpf unload
\end{lstlisting}
See \file{bpfMakefile}.

\subsection{Run the experiments and build this report}
\begin{lstlisting}[language=bash]
# build + load + all tests + plots
sudo ./run_all.sh -n 10 -t 10 -i 0.2 -r 5 --seed 42
# re-make plots and tables only
python3 run_experiments.py --plot-only results/<run>
# this PDF -> report/report.pdf (default: newest results/<run>)
python3 report/make_report.py [results/<run>]
\end{lstlisting}
\file{runallsh} builds and registers the program, runs \file{runexperimentspy} for all four
algorithms, and prints the program counters. \file{reportmakereportpy} (\pyReport) reads
\texttt{metadata.json}, \texttt{test\_order.csv}, \texttt{destinations.csv} and
\texttt{tables/*.csv} from the run directory and writes the tables and numbers in this report into
\texttt{report/generated/}. It finds the line ranges of the linked functions in the source, then runs
\texttt{pdflatex} on \file{reportreporttex}. Building the report needs only Python~3 and a TeX
distribution with \texttt{booktabs}, \texttt{listings}, \texttt{ulem} and \texttt{hyperref}.
A local network for testing without public servers is in \file{toolslocaltestbedsh}. The
run \texttt{results/20261018-133227} is a sanity check on that local network and is not used in this report.

\subsection{Raw data and configuration}
Everything is in \texttt{\rundir/}:
\begin{center}\small
\begin{tabular}{>{\raggedright\arraybackslash}p{5.4cm}>{\raggedright\arraybackslash}p{10.0cm}}
\toprule
File & Content \\
\midrule
\texttt{samples/}\allowbreak\texttt{r<rep>\_d<dest>\_<alg>.csv} & raw per-interval samples (all \texttt{TCP\_INFO} fields listed above) \\
\texttt{runs.jsonl} & one record per test: requested and actual algorithm, status, server results, errors \\
\texttt{metadata.json} & kernel, available/allowed algorithms, arguments, seed, \texttt{TCP\_INFO} layout, eBPF registration and counters before and after \\
\texttt{test\_order.csv} & exact execution order, attempts and status \\
\texttt{destinations.csv}, \texttt{selection\_log.csv} & chosen destinations and the selection probes \\
\texttt{tables/*.\{csv,md\}} & per-run, per-destination and aggregate tables \\
\texttt{plots/*.pdf} & every figure in this report \\
\bottomrule
\end{tabular}
\end{center}
Dependencies are in \texttt{requirements.txt} (numpy, matplotlib). The server list is \texttt{servers.csv}
(from iperf3serverlist.net, parsed by \file{serverspy}).

\subsection{Where the code is}
\begin{center}\small
\begin{tabular}{>{\raggedright\arraybackslash}p{6.2cm}>{\raggedright\arraybackslash}p{9.2cm}}
\toprule
What & Code \\
\midrule
Parameters / per-connection state & \grParams, \grState \\
Initial state & \grInit \\
Loss response ($\beta=0.7$) & \grSsthresh \\
Additive increase / slow start & \grCongAvoid \\
RTT and throughput slow-start exit & \grPktsAcked \\
RTO re-arm / spurious-loss undo & \grSetState, \grUndo \\
Registration (\texttt{tcp\_congestion\_ops}) & \grOps \\
Select algorithm on data socket & \pyConnect, \pySetCC \\
Record actual algorithm & \pyGetCC, \pyRun \\
Sampling (cwnd, RTT, goodput, loss) & \pySample \\
Experiment loop and test order & \pyRunExperiments \\
eBPF counter snapshot & \pyBpfSnapshot \\
Tables and aggregate statistics & \pyBuildTables, \pyRunStats \\
Representative destination & \pyChooseRep \\
Figures \ref{fig:ts}, \ref{fig:sc}, \ref{fig:regimes} and Appendix~\ref{app:plots} & \pyMakePlots, \pyPlotTs, \pyPlotScatter, \pyPlotRegimes \\
This report & \pyReport \\
\bottomrule
\end{tabular}
\end{center}

% ===========================================================================
\section{Use of AI tools}
\aiuse

% ===========================================================================
\appendix
\section{Per-algorithm plots at the representative destination}
\label{app:plots}
These are the Assignment~2 Q2(b) figures made separately for each algorithm, at
\texttt{\meta{repdest}}, with the same axis limits for all four algorithms.

\newcommand{\algfig}[3]{% file stem, label, file stem with escaped underscore
\begin{figure}[H]
  \centering
  \includegraphics[height=0.56\textheight,keepaspectratio]{\plotdir q2_timeseries_#1.pdf}\\[4pt]
  \includegraphics[width=\textwidth,height=0.3\textheight,keepaspectratio]{\plotdir q2_scatter_#1.pdf}
  \caption{#2 (\texttt{#3}) at \texttt{\meta{repdest}}: time series (top, repetition 1) and scatter (bottom, all repetitions).
  Files: \texttt{plots/q2\_timeseries\_#3.pdf}, \texttt{plots/q2\_scatter\_#3.pdf}.}
\end{figure}}
\algfig{cubic}{CUBIC}{cubic}
\algfig{reno}{Reno}{reno}
\algfig{bbr}{BBR}{bbr}
\algfig{gentle_reno}{\GR{}}{gentle\_reno}

\begin{figure}[H]
  \centering
  \includegraphics[width=\textwidth]{\plotdir q3_timeseries_overlay.pdf}
  \caption{The four algorithms overlaid at \texttt{\meta{repdest}}, repetition 1. File: \texttt{plots/q3\_timeseries\_overlay.pdf}.}
\end{figure}

\begin{figure}[H]
  \centering
  \includegraphics[width=\textwidth,page=1]{\plotdir q1_goodput_timeseries.pdf}
  \caption{Goodput time series for every destination and algorithm (page 1 of 2). File: \texttt{plots/q1\_goodput\_timeseries.pdf}.}
\end{figure}
\begin{figure}[H]
  \centering
  \includegraphics[width=\textwidth,page=2]{\plotdir q1_goodput_timeseries.pdf}
  \caption{Goodput time series for every destination and algorithm (page 2 of 2).}
\end{figure}

\section{Test order}
\label{app:order}
Execution order of all \stat{order}{total} tests (\sout{struck out} = failed). Each line runs the
four algorithms back to back on one destination.
{\small\begin{longtable}{rrll}
\toprule
Rep & Dest & Destination & Algorithm order \\
\midrule
\endhead
1 & 1 & \texttt{41.\allowbreak{}110.\allowbreak{}39.\allowbreak{}130:5201} & CUBIC $\to$ Reno $\to$ BBR $\to$ Gentle Reno \\
1 & 2 & \texttt{akl.\allowbreak{}linetest.\allowbreak{}nz:5301} & Reno $\to$ BBR $\to$ Gentle Reno $\to$ CUBIC \\
1 & 3 & \texttt{lon.\allowbreak{}speedtest.\allowbreak{}clouvider.\allowbreak{}net:5207} & BBR $\to$ Gentle Reno $\to$ CUBIC $\to$ Reno \\
1 & 4 & \texttt{speedtest.\allowbreak{}sin1.\allowbreak{}sg.\allowbreak{}leaseweb.\allowbreak{}net:5207} & Gentle Reno $\to$ CUBIC $\to$ Reno $\to$ BBR \\
1 & 5 & \texttt{speedtest.\allowbreak{}wdc2.\allowbreak{}us.\allowbreak{}leaseweb.\allowbreak{}net:5203} & CUBIC $\to$ Reno $\to$ BBR $\to$ Gentle Reno \\
1 & 6 & \texttt{72.\allowbreak{}1.\allowbreak{}86.\allowbreak{}46:30006} & Reno $\to$ BBR $\to$ Gentle Reno $\to$ CUBIC \\
1 & 7 & \texttt{speed1.\allowbreak{}fiberby.\allowbreak{}dk:9212} & BBR $\to$ \sout{Gentle Reno} $\to$ CUBIC $\to$ \sout{Reno} \\
1 & 8 & \texttt{41.\allowbreak{}210.\allowbreak{}185.\allowbreak{}162:5201} & Gentle Reno $\to$ CUBIC $\to$ Reno $\to$ BBR \\
1 & 9 & \texttt{speedtest.\allowbreak{}nyc1.\allowbreak{}us.\allowbreak{}leaseweb.\allowbreak{}net:5207} & CUBIC $\to$ Reno $\to$ BBR $\to$ Gentle Reno \\
1 & 10 & \texttt{speedtest.\allowbreak{}phx1.\allowbreak{}us.\allowbreak{}leaseweb.\allowbreak{}net:5206} & Reno $\to$ BBR $\to$ Gentle Reno $\to$ CUBIC \\
2 & 1 & \texttt{41.\allowbreak{}110.\allowbreak{}39.\allowbreak{}130:5201} & Reno $\to$ BBR $\to$ Gentle Reno $\to$ CUBIC \\
2 & 2 & \texttt{akl.\allowbreak{}linetest.\allowbreak{}nz:5301} & BBR $\to$ Gentle Reno $\to$ CUBIC $\to$ Reno \\
2 & 3 & \texttt{lon.\allowbreak{}speedtest.\allowbreak{}clouvider.\allowbreak{}net:5207} & Gentle Reno $\to$ CUBIC $\to$ Reno $\to$ BBR \\
2 & 4 & \texttt{speedtest.\allowbreak{}sin1.\allowbreak{}sg.\allowbreak{}leaseweb.\allowbreak{}net:5207} & CUBIC $\to$ Reno $\to$ BBR $\to$ Gentle Reno \\
2 & 5 & \texttt{speedtest.\allowbreak{}wdc2.\allowbreak{}us.\allowbreak{}leaseweb.\allowbreak{}net:5203} & Reno $\to$ BBR $\to$ Gentle Reno $\to$ CUBIC \\
2 & 6 & \texttt{72.\allowbreak{}1.\allowbreak{}86.\allowbreak{}46:30006} & BBR $\to$ Gentle Reno $\to$ CUBIC $\to$ Reno \\
2 & 7 & \texttt{speed1.\allowbreak{}fiberby.\allowbreak{}dk:9212} & Gentle Reno $\to$ \sout{CUBIC} $\to$ Reno $\to$ \sout{BBR} \\
2 & 8 & \texttt{41.\allowbreak{}210.\allowbreak{}185.\allowbreak{}162:5201} & CUBIC $\to$ Reno $\to$ BBR $\to$ Gentle Reno \\
2 & 9 & \texttt{speedtest.\allowbreak{}nyc1.\allowbreak{}us.\allowbreak{}leaseweb.\allowbreak{}net:5207} & Reno $\to$ BBR $\to$ Gentle Reno $\to$ CUBIC \\
2 & 10 & \texttt{speedtest.\allowbreak{}phx1.\allowbreak{}us.\allowbreak{}leaseweb.\allowbreak{}net:5206} & BBR $\to$ Gentle Reno $\to$ CUBIC $\to$ Reno \\
3 & 1 & \texttt{41.\allowbreak{}110.\allowbreak{}39.\allowbreak{}130:5201} & BBR $\to$ Gentle Reno $\to$ CUBIC $\to$ Reno \\
3 & 2 & \texttt{akl.\allowbreak{}linetest.\allowbreak{}nz:5301} & Gentle Reno $\to$ CUBIC $\to$ Reno $\to$ BBR \\
3 & 3 & \texttt{lon.\allowbreak{}speedtest.\allowbreak{}clouvider.\allowbreak{}net:5207} & CUBIC $\to$ Reno $\to$ BBR $\to$ Gentle Reno \\
3 & 4 & \texttt{speedtest.\allowbreak{}sin1.\allowbreak{}sg.\allowbreak{}leaseweb.\allowbreak{}net:5207} & Reno $\to$ BBR $\to$ Gentle Reno $\to$ CUBIC \\
3 & 5 & \texttt{speedtest.\allowbreak{}wdc2.\allowbreak{}us.\allowbreak{}leaseweb.\allowbreak{}net:5203} & BBR $\to$ Gentle Reno $\to$ CUBIC $\to$ Reno \\
3 & 6 & \texttt{72.\allowbreak{}1.\allowbreak{}86.\allowbreak{}46:30006} & Gentle Reno $\to$ CUBIC $\to$ Reno $\to$ BBR \\
3 & 7 & \texttt{speed1.\allowbreak{}fiberby.\allowbreak{}dk:9212} & CUBIC $\to$ \sout{Reno} $\to$ BBR $\to$ \sout{Gentle Reno} \\
3 & 8 & \texttt{41.\allowbreak{}210.\allowbreak{}185.\allowbreak{}162:5201} & Reno $\to$ BBR $\to$ Gentle Reno $\to$ CUBIC \\
3 & 9 & \texttt{speedtest.\allowbreak{}nyc1.\allowbreak{}us.\allowbreak{}leaseweb.\allowbreak{}net:5207} & BBR $\to$ Gentle Reno $\to$ CUBIC $\to$ Reno \\
3 & 10 & \texttt{speedtest.\allowbreak{}phx1.\allowbreak{}us.\allowbreak{}leaseweb.\allowbreak{}net:5206} & Gentle Reno $\to$ CUBIC $\to$ Reno $\to$ BBR \\
4 & 1 & \texttt{41.\allowbreak{}110.\allowbreak{}39.\allowbreak{}130:5201} & Gentle Reno $\to$ CUBIC $\to$ Reno $\to$ BBR \\
4 & 2 & \texttt{akl.\allowbreak{}linetest.\allowbreak{}nz:5301} & CUBIC $\to$ Reno $\to$ BBR $\to$ Gentle Reno \\
4 & 3 & \texttt{lon.\allowbreak{}speedtest.\allowbreak{}clouvider.\allowbreak{}net:5207} & Reno $\to$ BBR $\to$ Gentle Reno $\to$ CUBIC \\
4 & 4 & \texttt{speedtest.\allowbreak{}sin1.\allowbreak{}sg.\allowbreak{}leaseweb.\allowbreak{}net:5207} & BBR $\to$ Gentle Reno $\to$ CUBIC $\to$ Reno \\
4 & 5 & \texttt{speedtest.\allowbreak{}wdc2.\allowbreak{}us.\allowbreak{}leaseweb.\allowbreak{}net:5203} & Gentle Reno $\to$ CUBIC $\to$ Reno $\to$ BBR \\
4 & 6 & \texttt{72.\allowbreak{}1.\allowbreak{}86.\allowbreak{}46:30006} & CUBIC $\to$ Reno $\to$ BBR $\to$ Gentle Reno \\
4 & 7 & \texttt{speed1.\allowbreak{}fiberby.\allowbreak{}dk:9212} & Reno $\to$ \sout{BBR} $\to$ Gentle Reno $\to$ \sout{CUBIC} \\
4 & 8 & \texttt{41.\allowbreak{}210.\allowbreak{}185.\allowbreak{}162:5201} & BBR $\to$ Gentle Reno $\to$ CUBIC $\to$ Reno \\
4 & 9 & \texttt{speedtest.\allowbreak{}nyc1.\allowbreak{}us.\allowbreak{}leaseweb.\allowbreak{}net:5207} & Gentle Reno $\to$ CUBIC $\to$ Reno $\to$ BBR \\
4 & 10 & \texttt{speedtest.\allowbreak{}phx1.\allowbreak{}us.\allowbreak{}leaseweb.\allowbreak{}net:5206} & CUBIC $\to$ Reno $\to$ BBR $\to$ Gentle Reno \\
5 & 1 & \texttt{41.\allowbreak{}110.\allowbreak{}39.\allowbreak{}130:5201} & CUBIC $\to$ Reno $\to$ BBR $\to$ Gentle Reno \\
5 & 2 & \texttt{akl.\allowbreak{}linetest.\allowbreak{}nz:5301} & Reno $\to$ BBR $\to$ Gentle Reno $\to$ CUBIC \\
5 & 3 & \texttt{lon.\allowbreak{}speedtest.\allowbreak{}clouvider.\allowbreak{}net:5207} & BBR $\to$ Gentle Reno $\to$ CUBIC $\to$ Reno \\
5 & 4 & \texttt{speedtest.\allowbreak{}sin1.\allowbreak{}sg.\allowbreak{}leaseweb.\allowbreak{}net:5207} & Gentle Reno $\to$ CUBIC $\to$ Reno $\to$ BBR \\
5 & 5 & \texttt{speedtest.\allowbreak{}wdc2.\allowbreak{}us.\allowbreak{}leaseweb.\allowbreak{}net:5203} & CUBIC $\to$ Reno $\to$ BBR $\to$ Gentle Reno \\
5 & 6 & \texttt{72.\allowbreak{}1.\allowbreak{}86.\allowbreak{}46:30006} & Reno $\to$ BBR $\to$ Gentle Reno $\to$ CUBIC \\
5 & 7 & \texttt{speed1.\allowbreak{}fiberby.\allowbreak{}dk:9212} & BBR $\to$ \sout{Gentle Reno} $\to$ CUBIC $\to$ \sout{Reno} \\
5 & 8 & \texttt{41.\allowbreak{}210.\allowbreak{}185.\allowbreak{}162:5201} & Gentle Reno $\to$ CUBIC $\to$ Reno $\to$ BBR \\
5 & 9 & \texttt{speedtest.\allowbreak{}nyc1.\allowbreak{}us.\allowbreak{}leaseweb.\allowbreak{}net:5207} & CUBIC $\to$ Reno $\to$ BBR $\to$ Gentle Reno \\
5 & 10 & \texttt{speedtest.\allowbreak{}phx1.\allowbreak{}us.\allowbreak{}leaseweb.\allowbreak{}net:5206} & Reno $\to$ BBR $\to$ Gentle Reno $\to$ CUBIC \\
\bottomrule
\end{longtable}
}

\end{document}
